% Geometría, topología y memoria. Se incluye desde manual_teorico_pedagogico.tex.

\section{Geometría, topología y memoria en sistemas dinámicos}
\label{sec:geometria-topologia-memoria}

\begin{objectivebox}
Construir una lectura universitaria del sistema dinámico como objeto geométrico:
identificar el espacio donde vive, pasar del problema de valor inicial al flujo,
reconocer invariancia y atracción, estudiar la deformación local de perturbaciones y
reducir órbitas periódicas a mapas de retorno. La última parte explica qué cambia
cuando la evolución conserva memoria mediante una derivada de Caputo y deriva el
esquema ABM--PECE desde la ecuación integral de Volterra.
\end{objectivebox}

Este capítulo reúne ideas que suelen estudiarse en cursos distintos. La topología
aporta el vocabulario de vecindad, frontera, continuidad y compacidad. La geometría
diferencial permite interpretar un campo como una dirección tangente en cada estado.
La teoría cualitativa de ecuaciones diferenciales organiza órbitas, conjuntos límite,
variedades invariantes y retornos. El análisis numérico convierte esas construcciones
en experimentos y conserva sus hipótesis explícitas. El propósito es aprender a moverse
entre esos cuatro niveles con sus funciones diferenciadas.

La secuencia de trabajo será siempre la misma. Primero se declara el espacio de
fases y el modelo; después se obtiene una propiedad analítica; luego se decide qué
gráfica puede hacerla visible; finalmente se delimita el alcance de la evidencia disponible.
Esta disciplina vincula cada figura con las conclusiones que su método sostiene.

\subsection{Topología del espacio de fases}

\paragraph{Abiertos, vecindades y continuidad.}
Una topología sobre un conjunto $X$ es una familia $\mathcal T$ de subconjuntos,
llamados abiertos, que contiene a $\varnothing$ y a $X$, es cerrada bajo uniones
arbitrarias y bajo intersecciones finitas. Una vecindad de $x$ es cualquier conjunto
que contiene un abierto al que pertenece $x$. Esta definición parece abstracta, pero
responde una pregunta concreta: ¿qué significa que dos estados sean próximos y que
una solución dependa continuamente de su condición inicial?

En $\R^n$ se usa normalmente la topología inducida por la distancia euclídea

\begin{equation}
 d_2(x,y)=\left(\sum_{j=1}^{n}|x_j-y_j|^2\right)^{1/2}.
 \label{eq:gt-metrica-euclidea}
\end{equation}

Una función $h:X\to Y$ es continua si la preimagen de todo abierto de $Y$ es abierta
en $X$. La formulación por preimágenes no depende de coordenadas particulares y por
eso resulta adecuada al comparar modelos escritos con variables distintas.

\paragraph{Interior, clausura y frontera.}
Para $A\subset X$, el interior $\operatorname{int}A$ reúne los puntos que poseen una
vecindad contenida en $A$; la clausura $\overline A$ añade todos sus puntos de
acumulación. La frontera es

\begin{equation}
 \partial A=\overline A\setminus\operatorname{int}A.
 \label{eq:gt-frontera}
\end{equation}

Toda vecindad de un punto de $\partial A$ encuentra puntos de $A$ y de su
complemento. Una frontera calculada sobre una malla sólo aproxima esta definición:
el último píxel de un color sólo registra una aproximación del conjunto matemático
$\partial A$.

\paragraph{Compacidad y conexidad.}
Un conjunto compacto de $\R^n$ es cerrado y acotado. Esta equivalencia no se
traslada automáticamente a espacios generales, pero en el curso será una herramienta
frecuente: permite extraer subsucesiones convergentes y evita que una órbita salga a
distancias arbitrarias dentro del conjunto considerado. Un conjunto conexo no puede
separarse en dos abiertos relativos, disjuntos y no vacíos. Compacidad controla la
posibilidad de tomar límites; conexidad controla la posibilidad de separar el objeto.

\paragraph{Producto y cociente.}
Algunos espacios de fases adoptan una geometría distinta de $\R^n$. En un péndulo, el ángulo satisface
$\theta\sim\theta+2\pi k$; por ello el estado natural es
$S^1\times\R$, no un rectángulo con dos bordes físicos diferentes. Una distancia
angular compatible es

\begin{equation}
 d_{S^1}(\theta_1,\theta_2)
 =\min_{k\in\mathbb Z}|\theta_1-\theta_2+2\pi k|.
 \label{eq:gt-distancia-circular}
\end{equation}

Si una curva cruza el borde derecho de una gráfica angular y reaparece en el
izquierdo, no ha saltado en el espacio de fases: la discontinuidad pertenece a la
carta usada para dibujarla. Esta distinción es decisiva al interpretar trayectorias,
retornos y distancias.

\begin{activitybox}{Práctica reproducible: una frontera exacta y su aproximación}
En \gui{Crear sistema} defina las variables \code{x,y,z}, sin parámetros, con
ecuaciones \code{x-x**3}, \code{-y} y \code{-z}. Use RK4,
$h=0.01$, $T=20$ y los estados iniciales
$(0.25,0.5,0.2)$, $(-0.25,0.5,0.2)$, $(0.01,0.5,0.2)$ y
$(-0.01,0.5,0.2)$. Guarde las series y los retratos con nombres que incluyan el
signo de $x_0$. Antes de simular, demuestre con una línea de fase que $x=0$ es la
frontera entre los destinos $x=1$ y $x=-1$. Después explique por qué cuatro
trayectorias ilustran esa conclusión, pero no la sustituyen.
\end{activitybox}

\paragraph{Lectura topológica de las gráficas.}
Una serie temporal informa cómo se recorre el estado; un retrato de fase informa qué
región se visita; una gráfica de condiciones iniciales informa qué destino se asignó
a cada punto. Ninguna vista contiene por sí sola las otras dos. Si una separación
parece brusca, hay que preguntar si procede de la dinámica, de la identificación de
coordenadas, de una tolerancia del clasificador o de la resolución elegida.

\begin{figure}[H]
\centering
\includegraphics[width=0.92\linewidth]{topology_neighborhood_boundary.png}
\caption{Lectura mínima de un conjunto en el espacio de fases: (a) una vecindad
abierta alrededor de un punto, (b) interior y frontera, y (c) clausura frente al
complemento. El trazo de la frontera es información geométrica; la pertenencia y la
apertura son propiedades del conjunto y de la topología elegida.}
\label{fig:gt-topologia-minima}
\end{figure}

\lecturaseccion{Eduardo Nahmad-Achar, \emph{Topología y geometría diferencial con
aplicaciones a la física}, capítulos sobre espacios topológicos, variedades y campos;
Morris W. Hirsch, Stephen Smale y Robert L. Devaney, \emph{Differential Equations,
Dynamical Systems, and an Introduction to Chaos}, capítulos 1--3.}

\subsection{Problema de valor inicial y flujo}

Considérese una ecuación autónoma

\begin{equation}
 \dot x=f(x),\qquad x(0)=x_0,\qquad x\in U\subset\R^n.
 \label{eq:gt-pvi}
\end{equation}

Si $f$ es localmente Lipschitz, el teorema de Picard--Lindelöf proporciona una
solución local única. La palabra \emph{local} importa: un campo polinómico puede
tener solución sólo hasta un tiempo finito si no se dispone de una cota que impida
el crecimiento sin límite. Cuando la solución existe en el intervalo requerido se
define

\begin{equation}
 \Phi_t(x_0)=x(t;x_0).
 \label{eq:gt-flujo-def}
\end{equation}

La autonomía y la unicidad producen la ley de composición

\begin{equation}
 \Phi_0=\operatorname{id},\qquad
 \Phi_{t+s}=\Phi_t\circ\Phi_s.
 \label{eq:gt-ley-flujo}
\end{equation}

La derivación es breve. Las funciones $u(t)=\Phi_{t+s}(x_0)$ y
$v(t)=\Phi_t(\Phi_s(x_0))$ resuelven la misma ecuación y comparten el dato inicial
$u(0)=v(0)=\Phi_s(x_0)$. Por unicidad, coinciden. Si cada solución puede prolongarse
hacia tiempos positivos y negativos, $\Phi_t$ es un flujo; si sólo se dispone de
$t\geq0$, se habla de semiflujo.

\paragraph{Órbita e invariancia.}
La órbita de $x_0$ es el conjunto de estados visitados,

\begin{equation}
 \mathcal O(x_0)=\{\Phi_t(x_0):t\in I_{x_0}\}.
 \label{eq:gt-orbita}
\end{equation}

Un conjunto $A$ es positivamente invariante si $\Phi_t(A)\subseteq A$ para todo
$t\geq0$. Es invariante si la igualdad se conserva para todo tiempo admisible. La
invariancia es una afirmación sobre todas las trayectorias que comienzan en $A$, no
una impresión obtenida de una curva dibujada dentro de él.

Para una región $N=\{x:h(x)\leq0\}$ con frontera regular, el signo de la derivada de
Lie

\begin{equation}
 L_fh(x)=\nabla h(x)^{\mathsf T}f(x)
 \label{eq:gt-derivada-lie}
\end{equation}

indica si el campo entra o sale. Si $L_fh<0$ en $\partial N$, el campo apunta
estrictamente hacia el interior y $N$ es positivamente invariante. Una integral
promedio negativa sobre la frontera requiere controlar el signo punto
por punto o usar un argumento global distinto.

\paragraph{Ejemplo resuelto: contracción anisótropa.}
Para

\begin{equation}
 \dot x=-x,\qquad \dot y=-2y,
 \label{eq:gt-lineal-contractivo}
\end{equation}

la solución es

\begin{equation}
 \Phi_t(x_0,y_0)=(e^{-t}x_0,e^{-2t}y_0).
 \label{eq:gt-flujo-contractivo}
\end{equation}

Cada eje y cada cuadrante cerrado son positivamente invariantes. Las órbitas se
aproximan al origen, pero no lo hacen a la misma velocidad: la componente $y$ pierde
una fracción fija más rápido. En un retrato de fase se obtiene
$y=Cx^2$ dentro de cada región de signo constante. La geometría de la curva elimina
el tiempo; la serie temporal conserva las tasas.

\begin{activitybox}{Ejercicio con pista: flujo e invariancia}
Para \eqref{eq:gt-lineal-contractivo}, compruebe directamente
\eqref{eq:gt-ley-flujo}. Después use
$V(x,y)=x^2+y^2$ y calcule $\dot V$. ¿Qué discos son positivamente invariantes?
\textbf{Pista:} $\dot V=-2x^2-4y^2$; compare el signo en cada circunferencia y no
confunda invariancia del disco con invariancia de cada circunferencia.
\end{activitybox}

\begin{figure}[H]
\centering
\includegraphics[width=0.90\linewidth]{flow_invariance_trapping.png}
\caption{Dos lecturas complementarias del flujo contractivo: (a) órbitas orientadas
por el campo y (b) transporte de una región $D$ hacia $\varphi_t(D)$. Que la imagen
quede contenida en $D$ expresa invariancia positiva de la región, no invariancia de
cada curva dibujada en su interior.}
\label{fig:gt-flujo-invariancia}
\end{figure}

\lecturaseccion{Jessica Angélica Jaurez Rosas, Laura Ortiz Bobadilla, Jesús Alberto
Palma Márquez y Ernesto Rosales González, \emph{Teoría geométrica de ecuaciones
diferenciales I}, capítulos sobre existencia, unicidad y flujo; Steven H. Strogatz,
\emph{Nonlinear Dynamics and Chaos}, capítulos 2 y 5.}

\subsection{Conjuntos límite, atracción y regiones atrapantes}

Una trayectoria larga puede continuar moviéndose aunque deje de visitar regiones
nuevas. El conjunto $\omega$-límite registra lo que permanece accesible a tiempos
arbitrariamente grandes:

\begin{equation}
 \omega(x)=\bigcap_{T\geq0}
 \overline{\{\Phi_t(x):t\geq T\}}.
 \label{eq:gt-omega}
\end{equation}

Si la semiórbita positiva está contenida en un compacto, $\omega(x)$ no es vacío,
es compacto y es positivamente invariante. El descarte de un transitorio en el
software aproxima la operación $t\geq T$, pero una ventana finita no ejecuta la
intersección para todos los $T$.

\paragraph{Atracción y cuenca.}
Sea $A$ un compacto invariante. Se dice que atrae a $x$ si

\begin{equation}
 \dist(\Phi_t(x),A)\longrightarrow0\qquad(t\to\infty).
 \label{eq:gt-atraccion}
\end{equation}

La cuenca correspondiente es

\begin{equation}
 \mathcal B(A)=\{x:\dist(\Phi_t(x),A)\to0\}.
 \label{eq:gt-cuenca}
\end{equation}

Esta definición separa el conjunto asintótico de los estados que llegan a él. Dos
atractores geométricamente parecidos pueden tener cuencas muy distintas; una cuenca
extensa tampoco vuelve más compleja la dinámica que ocurre dentro del atractor.

\paragraph{Región atrapante.}
Un compacto $N$ es atrapante si existe $t_*>0$ tal que
$\Phi_{t_*}(N)\subset\operatorname{int}N$. El conjunto

\begin{equation}
 A_N=\bigcap_{t\geq0}\Phi_t(N)
 \label{eq:gt-maximo-invariante}
\end{equation}

es el máximo conjunto invariante que permanece dentro de $N$, siempre que el flujo y
la compacidad permitan esta construcción. La región atrapante brinda una cota global;
requiere diagnósticos adicionales para clasificar $A_N$ como equilibrio, ciclo o conjunto más
complicado.

\paragraph{Ejemplo resuelto: oscilación radial autosostenida.}
Considérese

\begin{equation}
 \dot x=x-y-x(x^2+y^2),\qquad
 \dot y=x+y-y(x^2+y^2).
 \label{eq:gt-normal-radial}
\end{equation}

Con $r^2=x^2+y^2$ y $\theta=\operatorname{atan2}(y,x)$,

\begin{equation}
 \dot r=r(1-r^2),\qquad \dot\theta=1.
 \label{eq:gt-radial-polar}
\end{equation}

Para $0<a<1<b$, en $r=a$ se tiene $\dot r>0$ y en $r=b$ se tiene
$\dot r<0$. El anillo $a\leq r\leq b$ es atrapante. Toda condición con $r_0>0$
se aproxima a la circunferencia $r=1$, mientras el origen permanece como equilibrio
inestable. El retrato muestra una espiral hacia una curva cerrada; la serie temporal
muestra una oscilación cuya amplitud se estabiliza; la gráfica de $r(t)$ revela de
forma más directa la atracción radial.

\begin{activitybox}{Práctica reproducible: región atrapante y observable adecuada}
Introduzca \eqref{eq:gt-normal-radial} en \gui{Crear sistema} con variables
\code{x,y}, RK4, $h=0.01$ y $T=25$. Simule desde $(0.2,0)$, $(0.8,0)$,
$(1.4,0)$ y $(2,0)$. Use \gui{Retratos 2D} y \gui{Series temporales}; exporte
los datos y añada la columna $r=\sqrt{x^2+y^2}$. Compare $r(t)$ con la predicción
de signos de \eqref{eq:gt-radial-polar}. Repita con $h/2$ y registre el máximo de
$|r(t)-1|$ en la última cuarta parte de la simulación.
\end{activitybox}

\paragraph{Lectura de una cuenca calculada.}
Cada celda representa una condición inicial y una clasificación obtenida con un
horizonte, un paso y tolerancias. Una franja de celdas sin clasificación estable puede
deberse a una frontera fina, a transitorios largos o a un criterio demasiado rígido.
Antes de describir geometría, conviene seleccionar puntos a ambos lados de la franja,
duplicar el horizonte y refinar la malla localmente.

\begin{figure}[H]
\centering
\includegraphics[width=0.90\linewidth]{omega_limit_trapping_region.png}
\caption{Región atrapante y conjunto $\omega$-límite del ejemplo radial: (a) una
órbita entra al anillo y se aproxima al ciclo $r=1$; (b) la distancia al conjunto
límite decrece sin que la trayectoria converja a un punto fijo. La segunda vista
evita confundir atracción hacia un conjunto con convergencia de cada coordenada.}
\label{fig:gt-omega-atrapante}
\end{figure}

\lecturaseccion{Kathleen T. Alligood, Tim D. Sauer y James A. Yorke,
\emph{Chaos: An Introduction to Dynamical Systems}, secciones sobre conjuntos
límite, cuencas y fronteras; Tamás Tél y Márton Gruiz, \emph{Chaotic Dynamics: An
Introduction Based on Classical Mechanics}, capítulos sobre dinámica transitoria y
cuencas.}

\subsection{Conjugación topológica y equivalencia orbital}

Dos ecuaciones escritas de forma distinta pueden representar la misma organización de
órbitas. La noción precisa depende de cuánto se desea conservar.

\paragraph{Conjugación.}
Los flujos $\Phi_t$ sobre $X$ y $\Psi_t$ sobre $Y$ son topológicamente conjugados si
existe un homeomorfismo $h:X\to Y$ tal que

\begin{equation}
 h\circ\Phi_t=\Psi_t\circ h
 \qquad\text{para todo }t.
 \label{eq:gt-conjugacion}
\end{equation}

La conjugación conserva el reloj: estados correspondientes se observan en el mismo
tiempo. Transporta equilibrios, periodicidad, invariancia y estructura de conjuntos
límite. No conserva longitudes, ángulos ni tasas expresadas en una métrica arbitraria.

Si $h$ es un difeomorfismo y $y=h(x)$, la regla de la cadena da el campo transformado

\begin{equation}
 g(y)=Dh(h^{-1}(y))f(h^{-1}(y)).
 \label{eq:gt-empuje-campo}
\end{equation}

El factor $Dh$ es esencial. Cambiar sólo el nombre de las variables no transporta el
campo.

\paragraph{Equivalencia orbital.}
Si se permite una reparametrización creciente del tiempo, basta que $h$ lleve órbitas
orientadas en órbitas orientadas. Multiplicar el campo por una función positiva
$\alpha(x)$,

\begin{equation}
 \dot x=\alpha(x)f(x),\qquad \alpha(x)>0,
 \label{eq:gt-reparametrizacion}
\end{equation}

conserva las curvas integrales orientadas, pero cambia la rapidez con que se recorren.
Por eso la equivalencia orbital no conserva periodos ni exponentes por unidad de
tiempo.

\paragraph{Ejemplo de conjugación lineal.}
Para $\dot x=-x$, la solución es $x(t)=e^{-t}x_0$. Tome
$h(x)=x^3$, homeomorfismo de $\R$. Entonces

\begin{equation}
 y(t)=h(x(t))=e^{-3t}y_0,\qquad \dot y=-3y.
 \label{eq:gt-conjugacion-cubica}
\end{equation}

Los dos flujos son conjugados por $h$: tienen un equilibrio globalmente atractivo y la
misma organización de semiórbitas. Sus tasas lineales son diferentes. El ejemplo
advierte que una tasa numérica se presenta con su métrica y sus unidades, en vez de como invariante topológico.

\begin{activitybox}{Ejercicio con pista: transporte de una región}
Sea $x=Tz+c$, con $T$ invertible. Demuestre que la bola
$\|x-x_e\|_2\leq r$ se transforma en
$(z-z_e)^{\mathsf T}T^{\mathsf T}T(z-z_e)\leq r^2$.
Explique por qué una malla cuadrada en las coordenadas $z$ no representa en general
una malla cuadrada en $x$. \textbf{Pista:} sustituya
$x-x_e=T(z-z_e)$ antes de elevar la norma al cuadrado.
\end{activitybox}

\paragraph{Interpretación de gráficas comparadas.}
Dos retratos conjugados pueden verse estirados, inclinados o curvados. Lo importante
es si conservan la correspondencia entre órbitas, equilibrios y conexiones. Si además
se comparan series temporales, hay que aclarar si el reloj se conservó. Superponer
curvas sin transportar coordenadas y tiempo mezcla dos preguntas distintas.

\begin{figure}[H]
\centering
\includegraphics[width=0.92\linewidth]{conjugacy_orbital_equivalence.png}
\caption{Correspondencias de órbitas: (a) una órbita de $\varphi_t$, (b) su
deformación por un homeomorfismo $h$ con las mismas marcas temporales y (c) la misma
órbita geométrica con marcas redistribuidas por $\theta(t)$. El último panel ilustra
por qué equivalencia orbital conserva el recorrido, pero no necesariamente el reloj.}
\label{fig:gt-conjugacion-orbital}
\end{figure}

\lecturaseccion{Morris W. Hirsch, Stephen Smale y Robert L. Devaney,
\emph{Differential Equations, Dynamical Systems, and an Introduction to Chaos},
capítulos sobre equivalencia y linealización; John Guckenheimer y Philip Holmes,
\emph{Nonlinear Oscillations, Dynamical Systems, and Bifurcations of Vector Fields},
capítulo 1.}

\subsection{Variación, cambio de volumen y variedades invariantes}

Una trayectoria describe la evolución de un estado. Para estudiar la evolución de
una perturbación se deriva el flujo respecto de la condición inicial. Si
$x(t)=\Phi_t(x_0)$ y $f$ es diferenciable, la matriz fundamental

\begin{equation}
 M(t)=D_{x_0}\Phi_t(x_0)
 \label{eq:gt-matriz-fundamental}
\end{equation}

satisface la ecuación variacional

\begin{equation}
 \dot M(t)=Df(x(t))M(t),\qquad M(0)=I.
 \label{eq:gt-variacional}
\end{equation}

Una perturbación pequeña se aproxima por $\delta x(t)=M(t)\delta x_0$. Los valores
singulares de $M(t)$ describen los ejes principales del elipsoide obtenido al deformar
una esfera infinitesimal. Los exponentes finitos se calculan a partir de sus tasas
logarítmicas, pero dependen del horizonte y del tratamiento numérico.

\paragraph{Fórmula de Liouville.}
La identidad de Jacobi para el determinante da

\begin{equation}
 \frac{\dd}{\dd t}\det M
 =\det M\,\operatorname{tr}(M^{-1}\dot M).
 \label{eq:gt-jacobi-det}
\end{equation}

Al sustituir \eqref{eq:gt-variacional},
$M^{-1}\dot M=M^{-1}Df(x(t))M$, cuya traza es
$\operatorname{tr}Df(x(t))=\nabla\cdot f(x(t))$. Por integración,

\begin{equation}
 \det M(t)=
 \exp\left(\int_0^t\nabla\cdot f(x(s))\dd s\right).
 \label{eq:gt-liouville}
\end{equation}

La divergencia mide expansión instantánea de volumen. Divergencia negativa implica
contracción de elementos de volumen a lo largo del flujo, pero no identifica por sí
sola el conjunto al que llegan las trayectorias ni prueba que una región particular
sea atrapante.

\paragraph{Equilibrio hiperbólico y variedades.}
Un equilibrio $e$ es hiperbólico si $Df(e)$ no tiene valores propios con parte real
cero. El espacio tangente se separa como

\begin{equation}
 \R^n=E^s\oplus E^u,
 \label{eq:gt-descomposicion-hiperbolica}
\end{equation}

donde $E^s$ corresponde a valores propios con parte real negativa y $E^u$ a parte
real positiva. El teorema de la variedad estable proporciona objetos no lineales
$W^s_{\mathrm{loc}}(e)$ y $W^u_{\mathrm{loc}}(e)$ tangentes a esos subespacios en
$e$. La linealización determina la tangencia local, no la forma global completa.

\paragraph{Ejemplo de variedad invariante curvada.}
Considérese

\begin{equation}
 \dot x=x,\qquad \dot y=-y+x^2.
 \label{eq:gt-silla-no-lineal}
\end{equation}

El origen tiene valores propios $1$ y $-1$. La recta $x=0$ es la variedad estable.
Para buscar la variedad inestable como $y=h(x)=ax^2$, se impone invariancia:

\begin{equation}
 h'(x)\dot x=-h(x)+x^2.
 \label{eq:gt-invariancia-grafica}
\end{equation}

Como $h'(x)\dot x=2ax^2$ y $-h(x)+x^2=(1-a)x^2$, se obtiene
$2a=1-a$ y, por tanto,

\begin{equation}
 W^u_{\mathrm{loc}}(0):\quad y=\frac{x^2}{3}.
 \label{eq:gt-variedad-inestable}
\end{equation}

El subespacio lineal $E^u$ es el eje $x$; la variedad verdadera se separa de él en
orden cuadrático. Un retrato con escala amplia puede hacerlas parecer iguales.

\begin{activitybox}{Práctica reproducible: tangencia local y curvatura}
Defina \eqref{eq:gt-silla-no-lineal} en \gui{Crear sistema}. Use RK4,
$h=0.002$, $T=6$ y condiciones $(0.02,0.02^2/3)$,
$(-0.02,0.02^2/3)$, $(0.02,0)$ y $(0,0.2)$. En \gui{Retratos 2D}, superponga
la parábola $y=x^2/3$ mediante datos exportados o puntos de referencia. Repita con
$h/2$. Describa qué trayectorias permanecen cerca de la parábola y por qué la
proximidad visual se complementa con la ecuación de invariancia.
\end{activitybox}

\begin{figure}[H]
\centering
\includegraphics[width=0.94\linewidth]{variational_volume_manifolds.png}
\caption{Tres objetos del sistema variacional: (a) transporte de una perturbación por
$D\varphi_t$, (b) cambio de un elemento de área gobernado por el determinante y
(c) direcciones estable e inestable de una silla. Una elipse o un paralelogramo son
aproximaciones lineales locales; las variedades globales pueden curvarse.}
\label{fig:gt-variacion-volumen-variedades}
\end{figure}

\lecturaseccion{Jessica Angélica Jaurez Rosas, Laura Ortiz Bobadilla, Jesús Alberto
Palma Márquez y Ernesto Rosales González, \emph{Teoría geométrica de ecuaciones
diferenciales II}, secciones sobre primera variación y variedades invariantes; Steven
H. Strogatz, \emph{Nonlinear Dynamics and Chaos}, capítulo 6.}

\subsection{Sección de Poincaré, monodromía y teoría de Floquet}

Una órbita periódica de periodo $T$ es una curva cerrada en el espacio de fases. Para
estudiar su estabilidad se elige una hipersuperficie transversal

\begin{equation}
 \Sigma=\{x:h(x)=0\},\qquad
 \nabla h(x)^{\mathsf T}f(x)\neq0
 \label{eq:gt-seccion}
\end{equation}

en los cruces considerados. El primer retorno lleva $z\in\Sigma$ al siguiente cruce:

\begin{equation}
 P(z)=\Phi_{\tau(z)}(z).
 \label{eq:gt-retorno}
\end{equation}

La órbita periódica se convierte en un punto fijo $P(p)=p$. Si los autovalores de
$DP(p)$ tienen módulo menor que uno, el ciclo es localmente atractivo en direcciones
transversales. Un conjunto de cruces apiñados sugiere un punto fijo; demostrarlo exige
controlar el evento, el tiempo de retorno y el error.

\paragraph{Monodromía y multiplicadores.}
Sobre una solución periódica $\gamma(t)$, la ecuación variacional tiene coeficientes
periódicos:

\begin{equation}
 \dot M(t)=Df(\gamma(t))M(t),\qquad M(0)=I.
 \label{eq:gt-variacional-periodica}
\end{equation}

La matriz $M(T)$ es la matriz de monodromía y sus autovalores son los multiplicadores
de Floquet. En un flujo autónomo aparece un multiplicador $1$ asociado al
desplazamiento a lo largo de la órbita. Los multiplicadores restantes miden el
comportamiento transversal y coinciden, bajo las identificaciones apropiadas, con los
autovalores de $DP(p)$.

En un sistema plano, Liouville permite obtener el multiplicador transversal sin
integrar toda la matriz:

\begin{equation}
 \mu_\perp=
 \exp\left(\int_0^T\nabla\cdot f(\gamma(t))\dd t\right).
 \label{eq:gt-multiplicador-planar}
\end{equation}

\paragraph{Ejemplo de multiplicadores de Floquet.}
En \eqref{eq:gt-normal-radial}, el ciclo $r=1$ tiene periodo $T=2\pi$. Si
$r=1+\eta$, entonces

\begin{equation}
 \dot\eta=\left.\frac{\dd}{\dd r}[r(1-r^2)]\right|_{r=1}\eta
 +O(\eta^2)=-2\eta+O(\eta^2).
 \label{eq:gt-variacion-radial}
\end{equation}

Después de una vuelta, $\eta(T)\simeq e^{-4\pi}\eta(0)$; por tanto,
$\mu_\perp=e^{-4\pi}$ y el ciclo es fuertemente atractivo en dirección radial.
La curva cerrada del retrato no muestra por sí sola ese número: hace falta comparar
retornos o resolver la variación.

\begin{activitybox}{Práctica reproducible: construir un retorno desde datos}
Use la trayectoria de \eqref{eq:gt-normal-radial} con $(x_0,y_0)=(1.4,0.1)$,
RK4, $h=0.005$ y $T=30$. Exporte el CSV. Detecte cruces de $y<0$ a $y\geq0$
con $x>0$ e interpole linealmente el valor de $x$ en cada cruce. Construya
$x_{k+1}$ contra $x_k$ y compare $|x_k-1|$ entre vueltas consecutivas. Repita con
$h/2$. Informe el cociente tardío y compárelo con $e^{-4\pi}$, señalando la
limitación de la interpolación de eventos.
\end{activitybox}

\paragraph{Lectura de la gráfica de retorno.}
La diagonal $P(z)=z$ permite localizar puntos fijos. La pendiente indica estabilidad
en un mapa escalar: $|P'(z^*)|<1$ atrae y $|P'(z^*)|>1$ repele. Una nube gruesa puede
proceder de dinámica compleja, de una sección mal definida, de cruces tangenciales o
de muestreo insuficiente. Antes de interpretarla, deben revisarse orientación,
transitorio y refinamiento del instante de cruce.

\begin{figure}[H]
\centering
\includegraphics[width=0.90\linewidth]{poincare_floquet_geometry.png}
\caption{Geometría de retorno y estabilidad orbital: (a) una sección transversal
$\Sigma$ reduce el ciclo $\gamma$ a iteraciones discretas; (b) los multiplicadores de
Floquet se leen respecto del círculo unidad. El multiplicador trivial $\mu_1=1$
corresponde al desplazamiento de fase y no decide la estabilidad transversal.}
\label{fig:gt-poincare-floquet}
\end{figure}

\lecturaseccion{Frank C. Hoppensteadt, \emph{Analysis and Simulation of Chaotic
Systems}, secciones sobre estabilidad orbital y mapas de retorno; John Guckenheimer y
Philip Holmes, \emph{Nonlinear Oscillations, Dynamical Systems, and Bifurcations of
Vector Fields}, capítulos 1--3.}

\subsection{Grado, puntos fijos y existencia}

Los métodos topológicos de existencia no buscan necesariamente una fórmula cerrada.
Transforman el problema en un cero o un punto fijo y usan continuidad, compacidad y
comportamiento en la frontera para demostrar que el objeto debe existir.

\paragraph{Contracción.}
En un espacio métrico completo, una aplicación $T$ es contractiva si existe
$0<L<1$ tal que

\begin{equation}
 d(Tx,Ty)\leq Ld(x,y).
 \label{eq:gt-contraccion}
\end{equation}

El teorema de Banach garantiza un único punto fijo y la convergencia de
$x_{n+1}=T(x_n)$. La conclusión combina existencia, unicidad y un algoritmo; si sólo
se sabe que $T$ es continua sobre una región compacta, esas tres propiedades no se
obtienen juntas.

\paragraph{Brouwer y grado.}
Toda aplicación continua de una bola cerrada de $\R^n$ en sí misma posee un punto
fijo. El grado de Brouwer refina esta idea. Para
$F:\overline\Omega\to\R^n$ y $0\notin F(\partial\Omega)$ se define un entero
$\deg(F,\Omega,0)$. Si es distinto de cero, $F$ tiene al menos un cero en
$\Omega$. Cuando los ceros son regulares y finitos,

\begin{equation}
 \deg(F,\Omega,0)=
 \sum_{x\in F^{-1}(0)\cap\Omega}
 \operatorname{sgn}\det DF(x).
 \label{eq:gt-grado-regular}
\end{equation}

El grado se conserva durante una homotopía $F_\lambda$ mientras ningún cero cruce la
frontera. Garantiza existencia, pero no decide estabilidad ni describe la cuenca.

\paragraph{Ejemplo de punto fijo por contracción.}
Para $F(x)=x-x^3$ en $\Omega=(-2,2)$, los ceros son $-1,0,1$ y
$F'(x)=1-3x^2$. La suma de signos es

\begin{equation}
 \deg(F,(-2,2),0)=(-1)+(+1)+(-1)=-1.
 \label{eq:gt-grado-cubico}
\end{equation}

El grado no nulo obliga a conservar al menos un cero bajo toda deformación admisible.
Para saber cuáles de esos equilibrios atraen se necesita el signo de la linealización
o un argumento no lineal adicional.

\paragraph{Disparo y residual de retorno.}
En el problema de frontera

\begin{equation}
 u''=G(t,u,u'),\qquad u(0)=a,\qquad u(T)=b,
 \label{eq:gt-problema-frontera}
\end{equation}

se elige la pendiente inicial $s=u'(0)$ y se define

\begin{equation}
 S(s)=u(T;s)-b.
 \label{eq:gt-residual-disparo}
\end{equation}

Un cero de $S$ resuelve el problema. Para una órbita periódica se usa de manera
análoga el residual $R(z)=P(z)-z$. Aplicar grado a $R$ requiere que el retorno exista
y sea continuo en la región; el cálculo de retornos debe acompañarse de la verificación explícita de esas hipótesis.

\begin{activitybox}{Ejercicios con pistas: tres formas de existencia}
\textbf{1.} Resuelva por disparo $u''+u=0$, $u(0)=0$,
$u(\pi/2)=1$. \textbf{Pista:} con pendiente $s$, $u(t;s)=s\sin t$.

\textbf{2.} Si $P:[a,b]\to[a,b]$ es continuo, pruebe que tiene un punto fijo.
\textbf{Pista:} estudie los signos de $P(x)-x$ en los extremos.

\textbf{3.} Calcule el grado de $x^3-x$ en intervalos pequeños alrededor de cada
raíz y verifique aditividad al reunirlos.
\end{activitybox}

\begin{figure}[H]
\centering
\includegraphics[width=0.82\linewidth]{degree_fixed_point.png}
\caption{Lectura geométrica del grado en el plano: el campo sobre la frontera rodea
una vez al origen cuando el cero está dentro de $D$ y no lo rodea cuando está fuera.
El dibujo orienta el cálculo; la conclusión de existencia requiere continuidad y
ausencia de ceros sobre la frontera.}
\label{fig:gt-grado-existencia}
\end{figure}

\lecturaseccion{Pablo Gustavo Amster, \emph{Métodos topológicos en el estudio de las
ecuaciones diferenciales no lineales}, capítulos sobre disparo, punto fijo y grado;
Jessica Angélica Jaurez Rosas et al., \emph{Teoría geométrica de ecuaciones
diferenciales II}, secciones sobre existencia y continuación.}

\subsection{Homoclinicidad, herradura y dinámica simbólica}

Una conexión homoclínica aparece cuando una órbita sale de un objeto hiperbólico por
su variedad inestable y regresa por su variedad estable. En un mapa, si $p$ es un
punto fijo hiperbólico, una intersección

\begin{equation}
 q\in W^s(p)\cap W^u(p),\qquad
 T_qW^s+T_qW^u=T_qX,
 \label{eq:gt-homoclinica-transversal}
\end{equation}

es transversal. La transversalidad impide que la intersección desaparezca ante una
perturbación pequeña. El teorema de Smale--Birkhoff establece que una intersección
homoclínica transversal de un punto periódico hiperbólico genera, para alguna
iteración, una dinámica tipo herradura sobre un conjunto invariante.

\paragraph{Estirar, contraer y plegar.}
La herradura toma una región, la estira en una dirección, la contrae en otra y la
pliega de modo que dos tiras regresan a la región inicial. Cada órbita que permanece
en el conjunto puede codificarse según la tira visitada:

\begin{equation}
 \ldots s_{-1}.s_0s_1s_2\ldots,\qquad s_k\in\{0,1\}.
 \label{eq:gt-codigo-simbolico}
\end{equation}

El corrimiento $\sigma$ elimina el símbolo presente y desplaza la secuencia. En una
métrica como

\begin{equation}
 d(s,r)=\sum_{k\in\mathbb Z}\frac{|s_k-r_k|}{2^{|k|+1}},
 \label{eq:gt-metrica-secuencias}
\end{equation}

dos secuencias son próximas si coinciden alrededor del presente. El corrimiento posee
órbitas periódicas densas, transitividad y sensibilidad. Esta construcción transforma
geometría de intersecciones en una descripción combinatoria.

\paragraph{Definiciones de caos.}
Para un mapa continuo sobre un conjunto invariante, la definición de Devaney combina
transitividad, densidad de puntos periódicos y sensibilidad. La entropía topológica
mide crecimiento de órbitas distinguibles. Un exponente de Lyapunov positivo mide
expansión promedio a lo largo de una trayectoria. Estas nociones se relacionan en
contextos precisos, y cada una conserva su ámbito de aplicación.

\paragraph{Evidencia requerida para interpretar una figura.}
Un pliegue visible sugiere buscar variedades e intersecciones. Una nube de retornos
sugiere estudiar un mapa iterado. Una banda ancha sugiere ausencia de una frecuencia
dominante simple. Ninguna de esas señales encierra por sí misma una intersección
transversal, una conjugación con el corrimiento o un valor exacto de entropía.

\begin{activitybox}{Práctica reproducible: itinerarios del mapa logístico}
En \gui{Bifurcación} seleccione \gui{Logistico [listo]}, use $x_0=0.2$ y estudie
primero $r=3.2$ y después $r=4$. Para cada valor conserve 200 iteraciones después de
descartar 1000. Exporte la sucesión y codifique $s_n=0$ si $x_n<1/2$ y $s_n=1$ si
$x_n\geq1/2$. Repita desde $x_0+10^{-8}$. Compare la longitud del prefijo común,
los bloques de símbolos y los puntos del retrato $x_{n+1}$ contra $x_n$. Esta
práctica estudia itinerarios finitos; no presenta la muestra como demostración de una
conjugación global.
\end{activitybox}

\begin{activitybox}{Ejercicio con pista: periodicidad simbólica}
Construya las secuencias binarias de periodo primo $1$, $2$ y $3$ bajo el corrimiento.
¿Cuántas palabras de longitud $n$ existen antes de identificar rotaciones?
\textbf{Pista:} hay $2^n$ palabras; para contar órbitas de periodo primo hay que
eliminar palabras de periodo divisor y después identificar desplazamientos cíclicos.
\end{activitybox}

\begin{figure}[H]
\centering
\includegraphics[width=0.92\linewidth]{horseshoe_symbolic_dynamics.png}
\caption{Esquema docente de la herradura: (a) dos franjas codificadas, (b)
estiramiento y contracción, y (c) plegamiento con un itinerario binario asociado.
La figura explica el mecanismo y la codificación; verificar una herradura en un
sistema concreto exige construir el mapa y demostrar las intersecciones necesarias.}
\label{fig:gt-herradura-simbolica}
\end{figure}

\lecturaseccion{Robert L. Devaney, \emph{An Introduction to Chaotic Dynamical
Systems}, capítulos sobre dinámica simbólica; Morris W. Hirsch, Stephen Smale y
Robert L. Devaney, \emph{Differential Equations, Dynamical Systems, and an
Introduction to Chaos}, capítulos 15--16; Kathleen T. Alligood, Tim D. Sauer y James
A. Yorke, \emph{Chaos: An Introduction to Dynamical Systems}, capítulos sobre
herraduras y órbitas homoclínicas.}

\subsection{Estados con memoria y derivadas fraccionarias}

En una ODE autónoma, conocer $x(t)$ basta para continuar la solución porque la ley de
flujo permite reiniciar desde ese estado. Una ecuación de Caputo conserva información
del intervalo recorrido desde el terminal inferior. Para $0<q<1$,

\paragraph{Función Gamma e integración fraccionaria.}
La función Gamma prolonga el factorial y fija la normalización de los operadores:

\begin{equation}
 \Gamma(z)=\int_0^\infty s^{z-1}e^{-s}\dd s,
 \qquad \Gamma(z+1)=z\Gamma(z).
 \label{eq:gt-gamma}
\end{equation}

Para $\alpha>0$, la integral izquierda de Riemann--Liouville con terminal inferior
$t_0$ es

\begin{equation}
 (I_{t_0}^{\alpha}g)(t)=\frac{1}{\Gamma(\alpha)}
 \int_{t_0}^{t}(t-s)^{\alpha-1}g(s)\dd s.
 \label{eq:gt-integral-rl}
\end{equation}

La identidad beta--Gamma permite probar, bajo condiciones de integrabilidad que
justifican intercambiar el orden de integración,

\begin{equation}
 I_{t_0}^{\alpha}I_{t_0}^{\beta}=I_{t_0}^{\alpha+\beta}.
 \label{eq:gt-semigrupo-orden}
\end{equation}

La composición ocurre respecto del orden del operador. Su interpretación difiere de la composición
temporal para el estado de un problema con memoria. Distinguir ambas afirmaciones
evita atribuir a la solución una propiedad que sólo pertenece a las integrales.

\paragraph{Definiciones de derivada fraccionaria.}
Si $m-1<q<m$, las derivadas de Riemann--Liouville y de Caputo ordenan de modo
distinto la derivación entera y la integración fraccionaria:

\begin{equation}
 {}^{\mathrm{RL}}D_{t_0}^{q}x
 =\frac{\dd^m}{\dd t^m}I_{t_0}^{m-q}x,
 \qquad
 {}^{\mathrm C}D_{t_0}^{q}x
 =I_{t_0}^{m-q}x^{(m)}.
 \label{eq:gt-rl-caputo}
\end{equation}

Para $0<q<1$, Caputo anula una constante, mientras

\begin{equation}
 {}^{\mathrm{RL}}D_{t_0}^{q}c
 =\frac{c(t-t_0)^{-q}}{\Gamma(1-q)}.
 \label{eq:gt-rl-constante}
\end{equation}

La forma de Grünwald--Letnikov aproxima el operador mediante diferencias de todos los
pasos anteriores:

\begin{equation}
 {}^{\mathrm{GL}}D_{t_0}^{q}x(t_n)
 \approx h^{-q}\sum_{j=0}^{n}(-1)^j
 \binom{q}{j}x(t_{n-j}),
 \qquad
 \binom{q}{j}=\frac{\Gamma(q+1)}{\Gamma(j+1)\Gamma(q-j+1)}.
 \label{eq:gt-gl-discreta}
\end{equation}

Estas expresiones describen operadores distintos. Deben declararse el
operador, el terminal inferior, la regularidad y la información inicial. Caputo es
conveniente cuando se formulan datos ordinarios $x(t_0)$; Riemann--Liouville puede
requerir datos expresados mediante integrales fraccionarias; Grünwald--Letnikov ofrece
una ruta discreta que vuelve explícito el costo de conservar la historia.

\begin{activitybox}{Ejercicio de lectura activa: una potencia como banco de prueba}
Para $x(t)=(t-t_0)^\nu$, con $\nu>0$, use la identidad beta--Gamma para obtener
la integral \eqref{eq:gt-integral-rl} y después las derivadas de
Riemann--Liouville y Caputo cuando existan. Verifique por separado el caso constante.
Compare la fórmula con \eqref{eq:gt-gl-discreta} para $q=0.6$, $h=0.1,0.05$ y
$t-t_0=1$. La entrega debe indicar dominio de los parámetros y error relativo, no
sólo valores decimales.
\end{activitybox}

\lecturaseccion{Shantanu Das, \emph{Kindergarten of Fractional Calculus}, capítulos
1--3, para Gamma, generalización, integración y las formas de Grünwald--Letnikov,
Riemann--Liouville y Caputo; Edmundo Capelas de Oliveira, \emph{Solved Exercises in
Fractional Calculus}, capítulos 2--5, para funciones especiales, Mittag--Leffler,
transformadas y comparación de operadores.}

\paragraph{Terminal inferior e inicialización.}
El valor $t_0$ marca el comienzo de la memoria
incluida en el modelo. Cambiarlo modifica el operador. Si el proceso ya tenía historia
antes de $t_0$, esa preparación debe incorporarse mediante un término de
inicialización o una historia declarada. Reiniciar el cálculo en un tiempo intermedio
con el transporte de esa información se define un problema equivalente. La pregunta correcta se formula como
únicamente cuál es $x(t_0)$, sino qué preparación representa el dato inicial y desde
qué instante se acumula la memoria.

\lecturaseccion{Shantanu Das, \emph{Kindergarten of Fractional Calculus}, capítulo 5,
sobre inicialización; Richard Herrmann, \emph{Fractional Calculus: An Introduction
for Physicists}, capítulo 8, sobre no localidad y memoria.}

Con este marco, la derivada de Caputo de orden $0<q<1$ se escribe

\begin{equation}
 {}^{\mathrm C}D_{0+}^{q}x(t)
 =\frac{1}{\Gamma(1-q)}
 \int_0^t(t-s)^{-q}\dot x(s)\dd s.
 \label{eq:gt-caputo}
\end{equation}

El problema

\begin{equation}
 {}^{\mathrm C}D_{0+}^{q}x(t)=f(x(t)),\qquad x(0)=x_0
 \label{eq:gt-pvi-caputo}
\end{equation}

es equivalente, bajo las hipótesis habituales, a la ecuación de Volterra

\begin{equation}
 x(t)=x_0+\frac{1}{\Gamma(q)}
 \int_0^t(t-s)^{q-1}f(x(s))\dd s.
 \label{eq:gt-volterra}
\end{equation}

La equivalencia se obtiene aplicando la integral fraccionaria de orden $q$ y usando
$I^q\,{}^{\mathrm C}D^qx=x-x_0$. La integral de Volterra muestra la memoria de forma
directa: cada valor pasado del campo participa en la continuación con un peso que
depende de la distancia temporal.

\paragraph{Perfil de memoria.}
Fije un tiempo $t$ y escriba el futuro como $t+\theta$. Al separar
\eqref{eq:gt-volterra} en $[0,t]$ y $[t,t+\theta]$, se obtiene

\begin{align}
 x(t+\theta)={}&m_t(\theta)
 +\frac{1}{\Gamma(q)}
 \int_0^\theta(\theta-u)^{q-1}f(x(t+u))\dd u,\label{eq:gt-volterra-desplazada}\\
 m_t(\theta)={}&x_0+\frac{1}{\Gamma(q)}
 \int_0^t(t+\theta-s)^{q-1}f(x(s))\dd s.\label{eq:gt-perfil-memoria}
\end{align}

El punto actual es sólo $m_t(0)=x(t)$. Reiniciar con el dato $x(t)$ sustituye la
función $m_t(\theta)$ por un término constante y, en general, produce otra
continuación. El estado adecuado puede representarse en un espacio de funciones, por
ejemplo $\mathfrak C=C([0,\infty),\R^n)$ con convergencia uniforme en intervalos
compactos. Sobre ese espacio puede recuperarse una ley de semiflujo para los perfiles,
aunque no para las aplicaciones puntuales obtenidas al conservar sólo $x(t)$.

\paragraph{Función de Mittag--Leffler y decaimiento fraccionario.}
La ecuación escalar

\begin{equation}
 {}^{\mathrm C}D_{0+}^{q}x=\lambda x,\qquad x(0)=x_0
 \label{eq:gt-caputo-lineal}
\end{equation}

tiene solución

\begin{equation}
 x(t)=x_0E_q(\lambda t^q),\qquad
 E_q(z)=\sum_{k=0}^{\infty}\frac{z^k}{\Gamma(qk+1)}.
 \label{eq:gt-mittag-leffler}
\end{equation}

Si la evolución puntual formara un semigrupo, tendría que cumplirse

\begin{equation}
 E_q(\lambda(t+s)^q)=E_q(\lambda t^q)E_q(\lambda s^q).
 \label{eq:gt-falsa-composicion}
\end{equation}

Al comparar el término lineal en $\lambda$, el lado izquierdo contiene
$(t+s)^q/\Gamma(q+1)$ y el derecho
$(t^q+s^q)/\Gamma(q+1)$. Para $0<q<1$ y $t,s>0$ no coinciden. Cuando $q=1$,
$E_1(z)=e^z$ y se recupera la composición ordinaria.

Para $\lambda=-a<0$, la solución describe relajación. En vez del decaimiento
exponencial $e^{-at}$ del caso entero, $E_q(-at^q)$ presenta una cola lenta. Una
comparación útil muestra primero la escala lineal, donde domina el transitorio, y
después una escala logarítmica, donde la cola se distingue con claridad. El orden
$q$ cambia simultáneamente la forma del kernel y la ley de relajación; su interpretación debe
interpretarse como un paso temporal.

\begin{activitybox}{Ejercicio con pista: continuidad de la memoria}
Tome $q=0.8$, $\lambda=-1$, $x_0=1$, $t=1$ y $s=1$. Evalúe
$E_q[-(t+s)^q]$ y $E_q(-t^q)E_q(-s^q)$ truncando la serie cuando el siguiente
término sea menor que $10^{-10}$. Repita con $q=1$.
\textbf{Pista:} calcule cada serie con precisión suficiente y compare tanto el valor
como la diferencia relativa; el caso $q=1$ sirve como control.
\end{activitybox}

\paragraph{Interpretación gráfica de la memoria.}
El kernel $K_q(r)=r^{q-1}/\Gamma(q)$ es singular integrable en $r=0$ y decae
lentamente. Una gráfica de $K_q$ para varios órdenes ayuda a comparar el peso de
retardos recientes y antiguos, y se complementa con la memoria completa de una trayectoria:
la contribución real es el producto del kernel con $f(x(s))$. Una proyección
$(x_i,x_j)$ puede cruzarse porque perfiles distintos comparten las mismas dos
coordenadas actuales; ese cruce no contradice unicidad en el espacio funcional.

\lecturaseccion{Kai Diethelm, \emph{The Analysis of Fractional Differential
Equations}, capítulos sobre problemas de valor inicial de Caputo y ecuaciones de
Volterra; Igor Podlubny, \emph{Fractional Differential Equations}, capítulos 1--4;
Richard Herrmann, \emph{Fractional Calculus: An Introduction for Physicists},
capítulos 2--3, 5--6 y 8, para funciones, derivadas, cálculo, oscilación y memoria;
Shantanu Das, \emph{Kindergarten of Fractional Calculus}, capítulos 6--7, para
relajación de Mittag--Leffler y ecuaciones diferenciales fraccionarias.}

\subsection{ABM--PECE de memoria completa}

La ecuación de Volterra permite derivar un integrador sin tratar la derivada
fraccionaria como una diferencia local. Sea $t_n=nh$ y
$F_j=f(t_j,x_j)$. En $t_{n+1}$,

\begin{equation}
 x(t_{n+1})=x_0+
 \frac{1}{\Gamma(q)}\sum_{j=0}^{n}
 \int_{t_j}^{t_{j+1}}
 (t_{n+1}-s)^{q-1}f(s,x(s))\dd s.
 \label{eq:gt-volterra-intervalos}
\end{equation}

\paragraph{Predictor.}
Al aproximar $f$ por la constante $F_j$ en cada subintervalo,

\begin{equation}
 \int_{t_j}^{t_{j+1}}(t_{n+1}-s)^{q-1}\dd s
 =\frac{h^q}{q}\left[(n+1-j)^q-(n-j)^q\right].
 \label{eq:gt-integral-peso-predictor}
\end{equation}

Como $q\Gamma(q)=\Gamma(q+1)$, el predictor es

\begin{equation}
 x_{n+1}^{\mathrm P}=x_0+
 \frac{h^q}{\Gamma(q+1)}
 \sum_{j=0}^{n}b_{j,n+1}F_j,\qquad
 b_{j,n+1}=(n+1-j)^q-(n-j)^q.
 \label{eq:gt-abm-predictor}
\end{equation}

\paragraph{Corrector.}
Al integrar una interpolación lineal del campo y evaluar el extremo nuevo con el
predictor,

\begin{equation}
 x_{n+1}=x_0+
 \frac{h^q}{\Gamma(q+2)}
 \left[f(t_{n+1},x_{n+1}^{\mathrm P})+
 \sum_{j=0}^{n}a_{j,n+1}F_j\right],
 \label{eq:gt-abm-corrector}
\end{equation}

donde

\begin{equation}
 a_{j,n+1}=
 \begin{cases}
 n^{q+1}-(n-q)(n+1)^q, & j=0,\\
 (n-j+2)^{q+1}+(n-j)^{q+1}-2(n-j+1)^{q+1},
 &1\leq j\leq n.
 \end{cases}
 \label{eq:gt-abm-pesos}
\end{equation}

PECE resume cuatro acciones: predecir, evaluar, corregir y evaluar otra vez. Las sumas
comienzan en $j=0$ en todos los pasos; quitar términos antiguos cambia el modelo
numérico y exige una aproximación de memoria corta con error declarado.

Para $q=1$, los pesos del predictor valen uno y el corrector recupera la regla
trapezoidal compuesta. Este límite es una prueba básica de implementación. También
deben compararse $h$, $h/2$ y $h/4$, mantener fijo el terminal inferior y registrar

\begin{equation}
 \mathcal N=(q,h,T,x_0,\text{método},\text{política de memoria}).
 \label{eq:gt-contrato-fraccionario}
\end{equation}

\begin{figure}[H]
\centering
\includegraphics[width=0.94\linewidth]{fractional_memory_kernel_pece.png}
\caption{Tres superficies de lectura para un problema de Caputo: (a) kernel de
memoria para varios órdenes, (b) relajación de Mittag--Leffler comparada con el caso
exponencial y (c) pesos históricos del predictor ABM. La altura de un peso representa una
contribución que debe multiplicarse por la evaluación histórica correspondiente y sumarse
desde el terminal inferior fijo.}
\label{fig:gt-memoria-abm}
\end{figure}

\begin{activitybox}{Laboratorio numérico: relajación escalar}
Implemente \eqref{eq:gt-abm-predictor}--\eqref{eq:gt-abm-pesos} para
${}^{\mathrm C}D^qx=-x$, $x(0)=1$, $q=0.8$, $T=5$ y
$h=0.05,0.025,0.0125$. Compare con $E_q(-t^q)$ y reporte el error máximo en una
malla común. Como control dentro de Toolbox Chaos, defina $\dot x=-x$ en
\gui{Crear sistema} y compare Euler, Heun y RK4 mediante \gui{Comparar metodos}:
esta ejecución corresponde al límite entero; el cálculo con memoria conserva el historial completo.
Conserve tabla, configuración y una gráfica de error; no entregue únicamente curvas
superpuestas.
\end{activitybox}

\begin{warningbox}
Un método que reinicia periódicamente desde el último valor de $x$ no implementa la
misma ecuación de Caputo con terminal inferior fijo. Una coincidencia visual durante
un intervalo corto no corrige esa diferencia. Antes de comparar dos soluciones deben
coincidir operador, orden, terminal inferior, historia, política de memoria, paso y
horizonte.
\end{warningbox}

\lecturaseccion{Kai Diethelm, \emph{The Analysis of Fractional Differential
Equations}, sección del método predictor--corrector; K. Diethelm, N. J. Ford y A. D.
Freed, ``A predictor-corrector approach for the numerical solution of fractional
differential equations'', \emph{Nonlinear Dynamics}, 2002; Igor Podlubny,
\emph{Fractional Differential Equations}, secciones numéricas.}

\subsection{Comprobación de hipótesis y resultados}

Los conceptos anteriores forman una cadena, no una colección de palabras. El espacio
de fases fija qué estados y vecindades son legítimos. La existencia y unicidad
permiten hablar de flujo. La invariancia y las regiones atrapantes organizan el
comportamiento global. La variación describe perturbaciones; Poincaré convierte un
ciclo en un punto fijo; Floquet mide su estabilidad transversal. El grado responde
preguntas de existencia sin resolver necesariamente el objeto. Las intersecciones de
variedades conectan geometría con dinámica simbólica. En Caputo, toda la cadena debe
levantarse a un estado que conserve memoria.

\begin{longtable}{>{\raggedright\arraybackslash}p{0.19\linewidth}
                  >{\raggedright\arraybackslash}p{0.34\linewidth}
                  >{\raggedright\arraybackslash}p{0.37\linewidth}}
\toprule
Objeto & Pregunta matemática & Comprobación antes de interpretar una gráfica\\
\midrule
\endhead
Espacio de fases & ¿Qué representa un estado y qué puntos se identifican? & Unidades,
periodicidad angular, dimensión y distancia elegida.\\
Flujo & ¿Existe una evolución única y componible? & Regularidad del campo, intervalo
de existencia y dependencia de la condición inicial.\\
Invariancia & ¿Las trayectorias permanecen en el conjunto? & Signo del campo en la
frontera o argumento exacto de conservación.\\
Atracción & ¿La distancia al conjunto tiende a cero? & Varias condiciones iniciales,
transitorio, horizonte y definición de distancia.\\
Variedad estable & ¿Qué estados convergen al objeto hiperbólico? & Espectro local,
tangencia y continuidad al ampliar la integración.\\
Retorno & ¿El siguiente cruce define una función? & Sección, orientación, exclusión de
tangencias e interpolación del evento.\\
Floquet & ¿Cómo cambia una perturbación después de un periodo? & Periodo, integración
variacional, multiplicador trivial y convergencia numérica.\\
Grado & ¿Debe existir un cero o punto fijo dentro de la región? & Ausencia de ceros en
la frontera y continuidad durante la homotopía.\\
Dinámica simbólica & ¿Existe una codificación compatible con la iteración? & Región
invariante, partición, cobertura y correspondencia de itinerarios.\\
Memoria Caputo & ¿Qué historia determina la continuación? & Orden, terminal inferior,
perfil, política de memoria y comparación de pasos.\\
\bottomrule
\end{longtable}

\begin{activitybox}{Proyecto breve: una afirmación, tres niveles de evidencia}
Elija uno de estos objetos: región atrapante de \eqref{eq:gt-normal-radial}, variedad
de \eqref{eq:gt-silla-no-lineal}, retorno del ciclo radial o itinerarios del mapa
logístico. Entregue una página con: definición; derivación analítica; experimento en
Toolbox Chaos; dos refinamientos; interpretación de cada gráfica; y una frase precisa
sobre el alcance que ya se estableció. Otra persona debe poder repetir el cálculo con
la información escrita, sin instrucciones orales.
\end{activitybox}

\paragraph{Problemas de cierre.}
\begin{enumerate}
  \item Demuestre que la imagen de un conjunto conexo bajo una aplicación continua es
  conexa. Explique cómo usaría el resultado para revisar una transformación de
  coordenadas.
  \item Para $\dot x=x(1-x)$, determine flujo implícito, equilibrios, cuencas y
  conjuntos $\omega$-límite. Dibuje primero la línea de fase.
  \item Construya una función $h$ cuya subnivel describa un disco y aplique
  \eqref{eq:gt-derivada-lie} al sistema radial.
  \item Verifique \eqref{eq:gt-liouville} para un sistema lineal diagonal y compare
  con el producto de sus soluciones escalares.
  \item En \eqref{eq:gt-silla-no-lineal}, busque el siguiente término de una gráfica
  invariante $h(x)=a_2x^2+a_3x^3$ y explique por qué $a_3=0$.
  \item Derive \eqref{eq:gt-multiplicador-planar} a partir de
  \eqref{eq:gt-liouville} y del multiplicador tangencial igual a uno.
  \item Para $P(x)=rx(1-x)$, determine la estabilidad de los puntos fijos y describa
  qué cambia cuando un multiplicador cruza $-1$.
  \item Justifique por qué un grado no nulo de $P-I$ demuestra existencia de un punto
  fijo, pero no su estabilidad.
  \item Compare recurrencia en el estado puntual con recurrencia del perfil
  $m_t(\theta)$ de \eqref{eq:gt-perfil-memoria}.
  \item Pruebe que los pesos $b_{j,n+1}$ de
  \eqref{eq:gt-abm-predictor} son positivos y explique su comportamiento para datos
  muy antiguos.
\end{enumerate}

\paragraph{Criterios de resolución para los problemas 1--10.}
En el problema 1 use la caracterización de conexidad por separaciones. En el 2 separe
variables y conserve los signos antes de integrar. En el 3 tome
$h=x^2+y^2-R^2$. En el 4 use
$M(t)=\operatorname{diag}(e^{\lambda_1t},\ldots,e^{\lambda_nt})$. En el 5 sustituya
la serie en \eqref{eq:gt-invariancia-grafica}. En el 6 separe la dirección tangente de
la transversal. En el 7 derive $P'(x^*)$. En el 8 recuerde que el grado no contiene
tasas temporales. En el 9 construya dos perfiles con el mismo valor en cero. En el 10
use que $u^q$ es creciente para $q>0$.

\begin{evidencebox}{Criterio de avance}
El capítulo está dominado cuando puede explicarse a partir de definiciones y teoremas por qué una
frontera de píxeles requiere una definición de frontera matemática; la linealización es
local; cómo un ciclo se convierte en punto fijo de retorno; qué aporta un grado no
nulo; qué hipótesis geométrica precede a una herradura; y por qué en Caputo el valor
actual requiere la historia necesaria para continuar.
\end{evidencebox}
